\section{Introduction}\label{sec:introduction}
The reciprocal chiral extension of the Tri-Octagon model supplies matching native, carrier and receiver accounts through two finite profiles~\cite{p1}. Its published same-trajectory theorem then connected initial supply from an empty carrier and receiver to later total native credit on one constant-parameter solution~\cite[Theorem 3.1]{p2}. That sign theorem left two quantitative questions distinct: what fraction of an actually transferred stock returns, and how much native storage can be extracted in the first place?

This paper presents the accepted accessible-transfer and projection-capacity investigations~\cite{at,pc}. The structural result applies to the general published preparation class at every fixed finite positive receiver coupling. The fixed real rank-two registration has a one-dimensional kernel. Projection onto that line produces a second conserved native quantity and a necessary extraction ceiling that can be strictly smaller than the previously published one. We prove its exact algebraic content, including feasibility and equality states, and distinguish those states from endpoints reached by the coupled dynamics.

The quantitative result has narrower hypotheses: the specified coherent launch family, empty external preparation and an explicit sufficient receiver-coupling threshold. On a fixed comparison-time window, an estimate for the actual solution proves positive transferred stock, its temporary concentration in the receiver, subsequent net native recovery and exact receiver charging and backflow. The recovered fraction can be close to one while the stock itself is small. The native-credit witness is imported from the published theorem; the receiver-power estimates are separate conclusions.

Section~\ref{sec:model} fixes the shared model without restricting its preparation. Section~\ref{sec:projection} proves the projection-capacity results. Sections~\ref{sec:transfer}--\ref{sec:proof} state and prove the finite transfer certificate. Section~\ref{sec:amount} compares amounts and ceilings while preserving the original normalization of accessible transfer. All statements concern exchange time. The intended geometric route remains a subsequent correspondence problem.
